\chapter{Antecedentes}\label{Chap:Antecedentes}
Los Sistemas de Levitación Magnética (SLM) al día de hoy tienen diversas aplicaciones: rodamientos magnéticos \cite{hung}, trenes de levitación magnética \cite{motoharu} y sistemas magnéticos para lanzamiento de plataformas espaciales \cite{kee}. El uso de los SLM están dados por las ventajas que presentan: bajo ruido, baja vibración, la posibilidad de operar en condiciones de vacío, en aplicaciones industriales donde es demandado el posicionamiento preciso (incluso a escala nanométrica). 

Los SLM dada su configuración atractiva o repulsiva pueden ser estables o inestables. Además de tener un comportamiento no lineal como consecuencia de la interacción de los campos. Adicionalmente el fenómeno de atascamiento-deslizamiento se presenta en la operación de estos dispositivos. Este problema surge principalmente por la no-linealidad de la fricción y la denominada zona muerta. Esta característica plantea un reto en el diseño de controladores de posición. Aunque se han propuesto varias estrategias para resolver o minimizar este problema, en la práctica solo se han reportado estrategias basadas en controladores conmutados y controladores clásicos con doble efecto integral, lo cual convierte al diseño de controladores de posición, para este tipo de procesos, en un reto de control. Estos controladores, si bien no garantizan que el error tienda a cero si lo reduce significativamente. Por otro lado, es importante mencionar que tampoco presentan un esfuerzo excesivo en la variable de control, como ocurre con las estrategias basadas en controladores con zona muerta inversa\cite{rubio}, Modos Deslizantes \cite{yang}\cite{yeh}\cite{bessa} o controladores predictivos\cite{galuppini}. La configuración atractiva del SLM, donde se presentan puntos de equilibrio inestables, no se ha analizado con el fenómeno atascamiento-deslizamiento. Asimismo encontrar un proceso que presente estas características tampoco es simple; sin embargo, este fenómeno se presenta en el problema de SLM en configuración inestable.

En el trabajo \emph{''Control conmutado para un sistema de levitación magnética con atascamiento-deslizamiento''} del 2014 \cite{alcantara}, se presentó un controlador en lazo cerrado para un SLM que presenta el fenómeno de atascamiento-delizamiento, basado en un esquema de control conmutado. Los resultados más relevantes fueron: se logró eliminar el atascamiento-deslizamiento, manteniendo un nivel bajo de error en estado estacionario y se obtuvo un controlador conmutado de menor complejidad que otras soluciones reportadas en la bibliografía.

En 2021 se presentó la tesis de maestría \emph{''Estrategia de Control para el Problema de Zona Muerta en Procesos Electromecánicos''} \cite{perez} donde se implementó un esquema de control lineal para reducir el fenómeno de la zona muerta en un motor de imanes permanentes de corriente directa. En este trabajo se diseñó un controlador con doble efecto integral lográndose un bajo error en estado estable y que permitía a la planta estar por muy poco tiempo en la zona muerta, sin considerar todas la dinámicas y características de la nolinealidad de la zona muerta.

Los puntos positivos de estos trabajos fueron, implementar controladores que lograron materializarse en tiempo real, de menor complejidad que otros reportados y que cumplieron con los criterios de estabilidad y robustez.

Los puntos negativos, que se consideraron en \cite{alcantara} fueron, generar un criterio de conmutación exacto, lo cual no es trivial, ya que requiere determinar con cierto grado de precisión el Ciclo Límite generado con el controlador PI y que, por lo tanto, constituye una dificultad que incrementa la complejidad del diseño,  el análisis  de estabilidad está completamente abierto y no se logró hacer cero el error en estado estable.

En el caso de \cite{perez}, el  error en estado estable es mayor al alcanzado en \cite{alcantara}. Sin embargo, el esquema de control es mucho más simple, la prueba de estabilidad es inmediata y, hubo una cuestión fundamental, no se requieren conocer los parámetros de la zona muerta.

%\includegraphics[scale=3]{marino.pdf}